%% 03b_ia_alineacion.tex — Metodología de fine-tuning y alineación del agente
%% de triaje. Estructura modular en tres fases (dataset, QLoRA, validación).
%%
%% Cifras del conjunto de datos: sections/_datos.tex (fuente:
%% data/curation.json). Hiperparámetros: ai_module/train_qlora.py.
%% Neutralidad de hardware: todo entorno de cómputo se describe en términos
%% metodológicos abstractos (memoria disponible, escalado de parámetros).

\section{Metodología de \emph{fine-tuning} y alineación del agente de triaje}
\label{sec:ia-metodologia}

El agente de triaje de la \cref{sec:agente} no es un modelo de propósito general
consultado mediante ingeniería de \emph{prompts}, sino un modelo de lenguaje
abierto (\emph{open-weight}) especializado en la taxonomía de ciberseguridad del
escáner. Su construcción se formaliza como un ciclo de entrenamiento modular en
\textbf{tres fases metodológicas} secuenciales ---preparación del conjunto de
datos, alineación y cuantización, y validación de seguridad---, cada una con un
artefacto verificable y un criterio de aceptación explícito. El objetivo de
diseño es doble y asimétrico: maximizar la supresión de falsos positivos
heurísticos minimizando, por encima de todo, el falso \emph{negativo} ---que un
juicio del modelo descarte una vulnerabilidad real.

Toda la metodología es agnóstica a la plataforma de ejecución. El
entrenamiento se realiza sobre infraestructura de aceleración de \emph{hardware}
estándar con un mecanismo de \textbf{escalado dinámico de parámetros} que
selecciona la capacidad del modelo base y la ruta de cuantización en función de
la memoria de cómputo disponible en el entorno, sin dependencias de una
arquitectura, un fabricante ni un modelo de acelerador concretos. Esta
abstracción permite reproducir el ciclo sobre entornos de cómputo heterogéneos
sin modificar la metodología.

\subsection{Fase 1: Preparación y saneamiento del conjunto de datos}
\label{sec:ia-dataset}

La primera fase transforma la telemetría cruda del escáner (una fila JSONL por
sonda, \cref{sec:telemetria}) en un conjunto de instrucciones supervisadas,
curado y balanceado. La curación se organiza como una tubería de cuatro etapas
---ingesta y enriquecimiento, limpieza, estructuración, y balanceo con partición
estratificada--- y produce \textbf{\IADatasetTriaje{}~instancias de triaje} que
cubren las \IAClasesTaxonomia{}~clases de vulnerabilidad de mayor incidencia:
inyección SQL (SQLi), inyección de comandos (\emph{command injection}),
\emph{cross-site scripting} (XSS) y \emph{path traversal}. La composición
final, deliberadamente equilibrada entre los \IAVeredictos{}~veredictos de la
taxonomía, se resume en la \cref{tab:ia-dataset}.

\begin{table}[t]
  \centering
  \caption{Composición del conjunto de datos de triaje tras el saneamiento
  (fuente: manifiesto de curación \code{data/curation.json}).}
  \label{tab:ia-dataset}
  \footnotesize
  \begin{tabular}{lr}
    \toprule
    Veredicto / partición & Instancias \\
    \midrule
    \textsc{true\_positive}   & \IATriajeTP{}  \\
    \textsc{false\_positive}  & \IATriajeFP{}  \\
    \textsc{uncertain}        & \IATriajeUNC{} \\
    \midrule
    Entrenamiento (\IATriajeSplit{})      & \IATriajeTrain{} \\
    Validación                            & \IATriajeVal{}   \\
    \midrule
    \textbf{Total}            & \textbf{\IADatasetTriaje{}} \\
    \bottomrule
  \end{tabular}
\end{table}

El saneamiento aplica tres salvaguardas contra el aprendizaje espurio, todas
registradas en el manifiesto de curación para su auditoría:

\begin{itemize}
  \item \textbf{Depuración semántica de parámetros.} El motor de semántica de
  parámetros (\code{param\_semantics.py}) identifica los puntos de inyección
  cuya función es de control de flujo ---parámetros de redirección o de destino
  de URL como \code{next}, \code{redirect\_uri} o \code{return\_to}. Sobre esos
  parámetros, la aplicación valida el valor antes de usarlo, de modo que
  cualquier divergencia de respuesta que no vaya acompañada de evidencia dura
  (una firma de error del intérprete, un oráculo temporal confirmado o una
  ejecución observada) se explica por esa validación y no por una inyección. El
  etiquetado fuerza tales casos a \textsc{false\_positive}, eliminando la clase
  de falso positivo heurístico ---\emph{la etiqueta es errónea mientras la
  evidencia parece plausible}--- que más contamina a un clasificador
  descendente.
  \item \textbf{Anonimización anti-fuga.} La entrada que ve el modelo expone
  únicamente un \emph{endpoint} y un nombre de parámetro genéricos: la ruta
  original ---que en el banco de referencia codifica la clase y el veredicto---
  se conserva solo en metadatos privados, evitando que el modelo aprenda un
  atajo hacia la etiqueta. La deduplicación opera sobre la clave
  \emph{observable} (entrada anonimizada más la latencia discretizada en
  cubetas), de modo que la variación de milisegundos no filtre una instancia
  casi idéntica a través de la frontera de entrenamiento y validación.
  \item \textbf{Balanceo y partición estratificada.} La clase mayoritaria se
  submuestrea hasta un equilibrio $1{:}1{:}1$ entre veredictos, y la partición
  \IATriajeSplit{} es estratificada para que la validación conserve las tres
  clases. Cada veredicto sintetizado se contrasta contra el mismo evaluador de
  evidencia que emplea la ruta real, descartando ---en lugar de mal etiquetar---
  todo caso cuya evidencia no discrimine su etiqueta pretendida.
\end{itemize}

\subsection{Fase 2: Alineación y cuantización (QLoRA)}
\label{sec:ia-qlora}

La segunda fase especializa un modelo base \emph{open-weight} sobre la taxonomía
mediante \emph{fine-tuning} eficiente en parámetros, evitando el coste de un
ajuste completo. Se combina la cuantización de la matriz de pesos base con la
adaptación de bajo rango (QLoRA~\cite{dettmers2023qlora}): el modelo base se
carga con cuantización de \textbf{\IACuantBits{}~bits} en el tipo de dato
\emph{NormalFloat} (NF4) con doble cuantización, y el gradiente actualiza
únicamente adaptadores de bajo rango (LoRA~\cite{hu2022lora}) sobre los pesos
congelados. Este esquema preserva la estabilidad numérica ---el cómputo se
realiza en precisión de punto flotante extendida mientras el almacenamiento se
mantiene cuantizado--- y reduce la huella de memoria en un factor suficiente
para que el ciclo completo quepa en el presupuesto de memoria seleccionado
dinámicamente por el entorno.

Los adaptadores se insertan sobre las \textbf{\IALoraModulos{}~proyecciones
lineales} del bloque \emph{transformer} ---las cuatro proyecciones de atención
(\emph{query}, \emph{key}, \emph{value}, \emph{output}) y las tres del bloque de
proyección \emph{feed-forward} (\emph{gate}, \emph{up}, \emph{down})--- es decir,
el conjunto completo de módulos objetivo, y no solo la atención. Los
hiperparámetros de la adaptación se fijaron para un corpus reducido y
especializado, priorizando la regularización sobre el rendimiento
(\cref{tab:ia-qlora}): un rango $r = \IALoraR{}$ con escala
$\alpha = \IALoraAlpha{}$ (razón estándar $\alpha = 2r$) y una tasa de
\emph{dropout} de \IALoraDropout{} sobre los adaptadores.

\begin{table}[t]
  \centering
  \caption{Configuración de la adaptación eficiente en parámetros (QLoRA).}
  \label{tab:ia-qlora}
  \footnotesize
  \setlength{\tabcolsep}{4pt}
  \begin{tabular}{ll}
    \toprule
    Componente & Valor \\
    \midrule
    Cuantización de pesos base   & NF4 \IACuantBits{}~bits + doble cuant. \\
    Rango de la adaptación ($r$) & \IALoraR{} \\
    Escala ($\alpha$)            & \IALoraAlpha{} \ ($\alpha = 2r$) \\
    \emph{Dropout} de LoRA       & \IALoraDropout{} \\
    Módulos objetivo             & \IALoraModulos{} (atención + MLP) \\
    Selección del modelo base    & escalado dinámico por memoria \\
    \bottomrule
  \end{tabular}
\end{table}

La alineación no se limita a enseñar la taxonomía: incluye un paso de
\textbf{restricción de dominio} que inyecta una minoría (aproximadamente el
\IARuidoCognitivo{}) de instancias de ruido conversacional ---peticiones ajenas
al análisis de seguridad--- cuya salida objetivo es un centinela de rechazo
fijo. Tras el ajuste, cualquier entrada que no sea un hallazgo de triaje o una
petición de síntesis de \emph{payloads} colapsa a ese centinela, que la capa de
inferencia interpreta como un veredicto \textsc{restricted} y descarta. El
efecto es un agente que conserva su competencia analítica dentro del dominio y
carece de la superficie de comportamiento de asistente general, reduciendo el
espacio de entradas capaces de inducir una respuesta fuera de la taxonomía.

\subsection{Fase 3: Validación y evaluación de seguridad}
\label{sec:ia-validacion}

La tercera fase somete al adaptador resultante a una evaluación de aceptación
independiente del entrenamiento, orientada a la mitigación de alucinaciones. Se
apoya en tres pilares:

\paragraph{Conjunto de validación disjunto (\emph{Golden Set})}
La evaluación se realiza sobre un conjunto curado a mano, \emph{disjunto} del
corpus de entrenamiento, de vulnerabilidades políglotas (verdaderos positivos)
y de escenarios adversarios de falso positivo ---páginas de bloqueo de WAF,
reflexión codificada como entidades, \emph{jitter} de latencia, respuestas
cacheadas o páginas de error independientes de la entrada. A diferencia de la
validación en línea durante el entrenamiento, esta evaluación se conduce con el
mismo cliente asíncrono que el escáner usa en producción, de modo que mide el
comportamiento desplegado y no el ajuste al corpus.

\paragraph{Decodificación estructurada}
Para que ninguna sección del análisis dependa de que el modelo respete el
formato por convención, la generación se restringe a un esquema tipado ---la
misma fuente única de verdad que define las etiquetas de entrenamiento y las
restricciones del servidor. Con enmascaramiento de \emph{logits} guiado por
gramática el modelo no puede emitir un token no conforme; cuando esa
restricción no está disponible, la salida se valida y repara contra el esquema
y, si resulta irreparable, se adopta el veredicto seguro \textsc{uncertain} en
lugar de propagar una respuesta malformada.

\paragraph{Métricas estrictas por clase}
Dado que la métrica de pérdida por entropía cruzada no informa sobre si el
modelo pasaría por alto una vulnerabilidad real, la evaluación reporta métricas
de seguridad explícitas por clase (\cref{tab:ia-metricas}). La cantidad
gobernante es la \textbf{tasa de falsos negativos} (FNR): la probabilidad de que
un verdadero positivo se clasifique como falso positivo, es decir, el coste
asimétrico que debe minimizarse. Se acompaña de la \emph{tasa de degradación}
(un verdadero positivo sobre el que el modelo se abstiene con \textsc{uncertain}
en vez de confirmarlo), de la \emph{cobertura de la clase falso positivo} (su
capacidad de suprimir ruido genuino) y de la \emph{exactitud de rechazo} sobre
las instancias de ruido conversacional. Los errores de inferencia
(agotamiento de tiempo, respuesta malformada tras reintentos) se contabilizan
por separado y nunca se pliegan en las métricas de clasificación, para que un
\emph{backend} inestable no pueda enmascararse como un buen FNR. El adaptador
solo se promueve a producción si supera los umbrales fijados sobre estas
métricas: la evaluación funciona como una \emph{puerta de aceptación}
(\emph{gate}), no como un mero informe.

\begin{table}[t]
  \centering
  \caption{Métricas de seguridad por clase de la puerta de evaluación. La FNR
  es la magnitud gobernante por su coste asimétrico.}
  \label{tab:ia-metricas}
  \footnotesize
  \setlength{\tabcolsep}{4pt}
  \begin{tabular}{ll}
    \toprule
    Métrica & Definición \\
    \midrule
    FNR                 & $P(\text{pred}{=}\textsc{fp} \mid \text{oro}{=}\textsc{tp})$ \\
    Degradación         & $P(\text{pred}{=}\textsc{unc} \mid \text{oro}{=}\textsc{tp})$ \\
    Cobertura \textsc{fp} & $P(\text{pred}{=}\textsc{fp} \mid \text{oro}{=}\textsc{fp})$ \\
    Exactitud de rechazo & fracción de ruido conversacional rechazado \\
    Exactitud (3 clases) & concordancia global, de referencia \\
    Tasa de error        & inferencias fallidas (aisladas, no clasif.) \\
    \bottomrule
  \end{tabular}
\end{table}
